\chapter{Camera Models}
\label{ch:camera}

\begin{objectives}
\begin{itemize}
  \item Derive the pinhole projection and the roles of the intrinsic parameters.
  \item Model and invert the radial-tangential lens distortion used by EuRoC.
  \item State the project's camera measurement model in normalized coordinates, with its noise.
  \item Derive the projection Jacobian and understand what a single view cannot observe.
  \item Understand camera-IMU extrinsics and time offsets.
\end{itemize}
\end{objectives}

% ------------------------------------------------------------------------------
\section{The pinhole model}

An ideal pinhole camera lets light through a single point, the optical center. A scene point with camera-frame
coordinates $p_c = (X, Y, Z)$, with $Z$ along the optical axis, projects through the center onto an image plane.
By similar triangles, its image on a virtual plane at unit distance in front of the center is
\begin{equation}
  \pi(p_c) = \begin{bmatrix} x \\ y \end{bmatrix} = \begin{bmatrix} X/Z \\ Y/Z \end{bmatrix}.
  \label{eq:pinhole}
\end{equation}
These are the \emph{normalized image coordinates}: dimensionless, independent of the lens, and equal to the
tangents of the ray's angles from the optical axis. The vector $\bar x = (x, y, 1)$ is the \emph{bearing} of the point,
defined up to scale.

\begin{figure}[t]
\centering
\begin{tikzpicture}[scale=1.0]
  \coordinate (O) at (0,0);
  \draw[frameaxis, plotblue] (O) -- (6.2,0) node[right] {$z_c$ (optical axis)};
  \draw[frameaxis, plotorange] (O) -- (0,-1.6) node[below] {$y_c$};
  \fill (O) circle (1.6pt) node[above left] {$C$};
  % image plane at z = 1 (drawn at x = 2)
  \draw[thick, primary] (2,1.8) -- (2,-1.4);
  \node[primary, font=\footnotesize, align=center] at (2,2.15) {normalized plane $Z = 1$};
  % point
  \coordinate (P) at (5.6,-1.25);
  \fill[plotgreen] (P) circle (2pt) node[right] {$p_c = (X, Y, Z)$};
  \draw[plotgreen, dashed] (O) -- (P);
  \coordinate (I) at (2,{-1.25*2/5.6});
  \fill[conv] (I) circle (1.8pt);
  \node[conv, font=\footnotesize, right] at (2.05,-0.18) {$y = Y/Z$};
  \draw[<->, black!60] (0,0.5) -- (2,0.5) node[midway, above, font=\footnotesize] {$1$};
  \draw[<->, black!60] (0,0.95) -- (5.6,0.95) node[midway, above, font=\footnotesize] {$Z$};
\end{tikzpicture}
\caption{Pinhole projection in the $y$-$z$ plane. The ray from the optical center $C$ to the point intersects the
plane $Z = 1$ at height $Y/Z$. All points along the ray project to the same image point, so depth is lost.}
\label{fig:pinhole}
\end{figure}

Pixels are obtained by scaling and shifting the normalized coordinates with the \emph{intrinsic matrix}:
\begin{equation}
  \begin{bmatrix} u \\ v \\ 1 \end{bmatrix} = K\begin{bmatrix} x \\ y \\ 1 \end{bmatrix}, \qquad
  K = \begin{bmatrix} f_u & s & c_u \\ 0 & f_v & c_v \\ 0 & 0 & 1 \end{bmatrix},
\end{equation}
with focal lengths $f_u, f_v$ in pixels, principal point $(c_u, c_v)$, and skew $s$ (zero for modern sensors). The focal
length in pixels is the physical focal length divided by the pixel pitch. EuRoC's \code{cam0} has
$f_u = 458.654$, $f_v = 457.296$, $c_u = 367.215$, $c_v = 248.375$ on a $752\times480$ sensor, a horizontal field of view of
about $79^\circ$.

% ------------------------------------------------------------------------------
\section{Lens distortion}

Real lenses bend rays, most strongly near the image edges. The \emph{radial-tangential} (Brown--Conrady) model
used by EuRoC maps ideal normalized coordinates $(x, y)$ to distorted ones $(x_d, y_d)$:
\begin{align}
  x_d &= x\,(1 + k_1r^2 + k_2r^4) + 2p_1xy + p_2(r^2 + 2x^2), \\
  y_d &= y\,(1 + k_1r^2 + k_2r^4) + p_1(r^2 + 2y^2) + 2p_2xy, \qquad r^2 = x^2 + y^2,
\end{align}
after which $K$ is applied to $(x_d, y_d)$. For \code{cam0}, $(k_1, k_2, p_1, p_2) = (-0.2834, 0.07396, 1.94\times10^{-4},
1.76\times10^{-5})$. The negative $k_1$ produces \emph{barrel} distortion (\cref{fig:distortion}). At the middle of the left and right image
edges the radial factor $1 + k_1r^2 + k_2r^4$ is about $0.79$, which moves points by about 100 pixels; at the corners it is
about $0.73$, a shift of about 160 pixels. Ignoring distortion is therefore not an option.

\begin{figure}[t]
  \centering
  \includegraphics[width=0.84\textwidth]{fig_distortion.pdf}
  \caption{EuRoC \code{cam0} lens distortion. A regular grid in normalized coordinates (gray) is imaged by the real
  lens as the curved blue grid. Distortion pulls points toward the center, increasingly toward the edges, so the
  sensor (orange) sees a wider field than an ideal pinhole with the same focal length would.}
  \label{fig:distortion}
\end{figure}

\paragraph{Undistortion.} The front end measures distorted pixels and needs ideal normalized coordinates. The model
has no closed-form inverse, but a fixed-point iteration converges quickly for moderate distortion: starting from
$x^{(0)} = x_d$,
\begin{equation}
  x^{(j+1)} = \frac{x_d - \Delta_x(x^{(j)}, y^{(j)})}{1 + k_1r_j^2 + k_2r_j^4},\qquad
  y^{(j+1)} = \frac{y_d - \Delta_y(x^{(j)}, y^{(j)})}{1 + k_1r_j^2 + k_2r_j^4},
\end{equation}
where $\Delta$ denotes the tangential terms. A Gauss--Newton iteration with the $2\times2$ Jacobian of the distortion
converges faster and more reliably near the edges. OpenCV's \code{cv::undistortPoints} implements the iterative
scheme. Milestone 8 checks undistortion by distorting and undistorting a grid of points.

\begin{keyidea}[undistort points, not images]
The estimator needs a few hundred feature locations, not a rectified image. Undistorting only the tracked points
costs microseconds; remapping a full image costs milliseconds, resamples the pixels (blurring them), and changes
the noise of the feature detector in a position-dependent way. Track in the raw image and undistort the tracks.
\end{keyidea}

Wide-angle and fisheye lenses need different models, such as the equidistant (Kannala--Brandt) model, but the
structure of the pipeline is the same: everything downstream of undistortion works with bearings.

% ------------------------------------------------------------------------------
\section{The measurement model}

\begin{convention}[camera model, AGENTS.md C6]
In normalized, undistorted image coordinates,
\begin{equation}
  p_c = R_{bc}\T\big(R_{wb}\T(p_f - p_{wb}) - p_{bc}\big), \qquad
  z = \begin{bmatrix} p_{c,x}/p_{c,z} \\ p_{c,y}/p_{c,z} \end{bmatrix} + n_{uv},
  \qquad \sigma_{uv} \approx \sigma_{px}/f .
\end{equation}
\end{convention}

The first equation chains the world-to-body and body-to-camera transformations of \eqref{eq:world-to-camera}. The noise
model follows from $u = f_u x + c_u$: a pixel error $\delta u$ corresponds to a normalized error $\delta u/f_u$. With
$\sigma_{px} = 1$~px and $f \approx 458$~px, $\sigma_{uv} \approx 2.2\times10^{-3}$. Strictly, distortion makes the mapping
from pixel noise to normalized noise position-dependent, by the inverse of the distortion Jacobian, and $f_u \neq f_v$.
Both effects are small for EuRoC and conventionally ignored, but they are a legitimate refinement.

% ------------------------------------------------------------------------------
\section{The projection Jacobian}

\begin{proposition}[Jacobian of the pinhole projection]
\label{prop:proj-jacobian}
\begin{equation}
  \frac{\partial\pi}{\partial p_c} = \frac{1}{Z}\begin{bmatrix} 1 & 0 & -X/Z \\ 0 & 1 & -Y/Z \end{bmatrix}
  = \frac1Z\begin{bmatrix} I_2 & -\pi(p_c) \end{bmatrix}.
\end{equation}
\end{proposition}
\begin{proof}
Differentiate $X/Z$ and $Y/Z$ with the quotient rule.
\end{proof}

The Jacobian has rank two, and its nullspace is spanned by $p_c$ itself: moving the point along its viewing ray does
not change its image. This is the single-view loss of depth in differential form, and it is why a feature must be
observed from at least two sufficiently different viewpoints before it constrains anything (\cref{ch:multiview}).
The factor $1/Z$ says that distant points have small image motion for a given change in position. They constrain
rotation well and translation poorly.

The full measurement Jacobian with respect to the filter state is the chain-rule product of this matrix with the
Jacobian of $p_c$ with respect to the clone pose and the feature position. Deriving those is part of
\cref{ch:msckfupdate}.

\begin{pitfall}[points behind the camera]
\eqref{eq:pinhole} happily projects points with $Z < 0$, flipping their image through the center. A triangulation or
linearization that produces a negative depth must be rejected (the \emph{cheirality} check), not projected. Points
with very small positive $Z$ produce huge Jacobians and should be rejected too.
\end{pitfall}

% ------------------------------------------------------------------------------
\section{Extrinsics and time offset}

\paragraph{Extrinsic calibration.} The transformation $T_{bc}$ between the IMU and the camera is estimated offline,
for EuRoC with Kalibr~\cite{furgale2013kalibr}, and stored as \code{T\_BS} in \code{cam0/sensor.yaml}. Its rotation is close to a
permutation of axes, and its translation is a few centimeters (\cref{app:euroc}). An error of $\delta\phi$ in the extrinsic
rotation biases every predicted bearing by about $\delta\phi$, which looks to the filter like a systematic residual. A
translation error matters mostly for near points. The project treats $T_{bc}$ as exact (AGENTS.md C3). Estimating it
online is a stretch goal.

\paragraph{Time offset.} If the camera timestamp differs from the IMU clock by $t_d$, a feature moving across the image at
velocity $\dot z$ is measured with an error of about $\dot z\,t_d$. A 1~ms offset at an image-plane velocity of 0.5 normalized
units per second (fast rotation) is an error of $5\times10^{-4}$, a quarter of a pixel. EuRoC is hardware synchronized, so
$t_d \approx 0$; consumer devices can have offsets of tens of milliseconds and must estimate it
online~\cite{li2014timeoffset}.

\paragraph{Shutter.} EuRoC's cameras have global shutters: every pixel is exposed at the same instant. Rolling-shutter
cameras expose rows sequentially, so each row has its own timestamp and pose, which requires an extended measurement
model.

\begin{takeaways}
\begin{itemize}
  \item Normalized coordinates $(X/Z, Y/Z)$ are lens-independent bearings; $K$ maps them to pixels.
  \item EuRoC's lens has strong barrel distortion; undistort the tracked points, not the images.
  \item The measurement model works in normalized coordinates with $\sigma_{uv} \approx \sigma_{px}/f$.
  \item $\partial\pi/\partial p_c = \frac1Z[I_2\ \ {-\pi(p_c)}]$; its nullspace is the viewing ray.
  \item Reject points with negative or tiny depth. Extrinsic and timing errors appear as systematic residuals.
\end{itemize}
\end{takeaways}

\begin{exercises}
  \item Compute the horizontal and vertical fields of view of \code{cam0} with and without distortion.
  \item Implement distortion and fixed-point undistortion for the EuRoC model. How many iterations are needed for
        $10^{-6}$ accuracy at the image corners? Compare with Gauss--Newton.
  \item Show that the normalized noise of a feature at the image corner, mapped through the inverse distortion
        Jacobian, differs from $\sigma_{px}/f$, and by how much.
  \item A point at depth 10~m is observed by a camera translating sideways at 1~m/s and rotating at 1~rad/s. Which motion
        dominates its image velocity? What does this imply for the observability of translation from distant points?
  \item Show that a point at infinity (direction $d$) has image $\pi(R\T d)$ independent of camera translation.
\end{exercises}
